\section{引言}\label{sec:introduction}
辐射源定位常先估计到达时间差等中间量，再由这些量求解位置。广义互相关与双曲线定位为这一流程提供了成熟工具\cite{KnappCarter1976,ChanHo1994}。当多个源使用重叠频带且波形未知时，各站接收的是不同传播响应下的混合信号，中间量的提取与跨站关联随之变难。直接位置估计（direct position determination，DPD）将传播关系直接带入观测模型，以位置作为主要待估参数，为多源联合处理提供了另一条途径\cite{ChenHudsonYao2002,WeissAmar2005}。

本文关注四个同步单天线站接收两个独立宽带源的情形。两源中心频率和带宽相同，接收功率可以明显不同，发射波形与发射功率未知。各站满足同一确定性传播规律，待解决的是如何从一次有限时长观测中高效求出两个位置。对这一问题，仅在单源谱中取两个峰容易把强源的邻峰当作第二个源；直接对两个二维位置做全网格组合搜索，又需要与网格点数平方成正比的计算。即使采用逐源求解，候选进入哪个局部区域、局部更新时是否重新估计波形，也会影响精度和成本。

宽带最大似然与交替投影（alternating projection，AP）通过固定部分位置、更新其余位置来降低联合搜索的维数\cite{ChenHudsonYao2002}。MUSIC型DPD利用噪声子空间检验位置响应\cite{WeissAmar2005}，MVDR借助协方差逆矩阵形成抑制干扰的空间评分\cite{Capon1969}；DPD-HR进一步研究了MVDR思想下的直接定位\cite{TirerWeiss2016}。这些准则依赖不同的观测与统计条件，比较时需要统一有效频带、传播信息和输入样本。

稀疏参数估计将粗搜索与离网精修结合起来。Newtonized orthogonal matching pursuit（NOMP）在贪心发现后进行Newton更新和循环精修\cite{Mamandipoor2016}。多频率二维块NOMP已用于离网波束形成\cite{Yang2022}；QNOMP采用准Newton联合优化改善多径参数估计\cite{Zeng2025}。这些研究已涵盖多频率建模、二维精修和联合优化。本文关注未知宽带波形怎样消元、强弱源候选怎样保留，以及这些选择如何影响实际定位成本。

可分离非线性最小二乘的变量投影思想，可将线性系数从非线性参数优化中消去\cite{GolubPereyra1973}。在本文场景中，位置为非线性参数，每个频点的源信号系数为线性参数。若局部位置试探仍使用固定的信号系数，所见曲率会包含本可由信号重估吸收的变化。另一方面，粗条件谱的排序并不保证精修后的排序。前者影响更新效率，后者影响可到达的解，因而需要分别处理。

围绕这些问题，本文形成以下三项工作。
\begin{enumerate}
  \item 针对同步单天线站的未知宽带强弱重叠源定位，建立统一的基带观测、功率与波形消元模型，将逐源发现、条件搜索及联合精修写在同一最小二乘目标下，明确不同求解策略的关系。
  \item 提出多候选联合精修AP：保留少量条件候选，联合调整源位置并重估波形，以精修后的目标选择分支。给出条件投影与联合拟合的等价关系，分析候选保留和联合更新分别承担的作用。
  \item 提出波形消元块NOMP：保留残差发现与循环更新，改用消元后的局部位置目标。通过Schur补刻画固定系数与消元曲率的差别，并结合受控对照和运算计数分析其计算优势。
\end{enumerate}

下文给出统一模型、两种求解算法及其计算关系，随后以基线比较、候选拆分和受控实验检验方法效果。扩大布局的比较沿用相同模型与评价口径，汇总结果经核验后纳入最终稿。
